\section{Introduction}

Gauss's work on the constructibility of regular polygons is one of his most famous achievements. At age 19, he proved that a regular heptadecagon \index{heptadecagon} (17-gon) is constructible \index{constructible polygon} with ruler and compass \index{ruler-and-compass construction}, and more generally characterized all constructible regular $n$-gons. This discovery was so important to Gauss that he requested a 17-gon be inscribed on his tombstone.

Gauss's diary entry from 1796 records his breakthrough: "EYPHKA! num = $\Delta$ $\Delta$ $\Delta$," referring to the construction of the regular 17-gon. This was the first progress on the polygon construction problem since the time of the ancient Greeks.

In this chapter, we cover:
\begin{itemize}
    \item Roots of unity and cyclotomic polynomials
    \item Gauss's criterion for constructible polygons
    \item The 17-gon construction using Gaussian periods \index{Gaussian period}
    \item General $n$-gon construction when possible
    \item Implementation of the construction
\end{itemize}

\section{Roots of Unity}

The vertices of a regular $n$-gon inscribed in the unit circle are the $n$-th roots of unity:
\[
\zeta_n^k = e^{2\pi i k / n}, \quad k = 0, 1, \ldots, n-1.
\]

\begin{definition}[Cyclotomic Polynomial]\index{cyclotomic polynomial}
The $n$-th cyclotomic polynomial is
\[
\Phi_n(x) = \prod_{\substack{1 \le k \le n \\ \gcd(k,n)=1}} (x - \zeta_n^k).
\]
It is the minimal polynomial of $\zeta_n$ over $\Q$, has degree $\varphi(n)$, and is irreducible.
\end{definition}

The cyclotomic polynomials satisfy the product formula:
\[
x^n - 1 = \prod_{d|n} \Phi_d(x).
\]
This allows recursive computation: $\Phi_n(x) = \frac{x^n - 1}{\prod_{d|n, d<n} \Phi_d(x)}$.

\begin{example}[First Cyclotomic Polynomials]
$\Phi_1(x) = x - 1$, $\Phi_2(x) = x + 1$, $\Phi_3(x) = x^2 + x + 1$, $\Phi_4(x) = x^2 + 1$, $\Phi_5(x) = x^4 + x^3 + x^2 + x + 1$.
\end{example}

\begin{theorem}[Euler's Totient Function]\index{Euler's totient function}
The degree of $\Phi_n(x)$ is $\varphi(n) = n\prod_{p|n}(1-1/p)$, where the product is over distinct prime divisors of $n$.
\end{theorem}

\section{Gauss's Criterion}

\begin{theorem}[Gauss-Wantzel Theorem]\index{Gauss-Wantzel theorem}
A regular $n$-gon is constructible with ruler and compass if and only if
\[
n = 2^k p_1 p_2 \cdots p_m
\]
where $p_i$ are distinct Fermat prime \index{Fermat prime}s: $F_k = 2^{2^k} + 1$.
\end{theorem}

Known Fermat primes: $3, 5, 17, 257, 65537$. The next $2^{2^5}+1 = 4294967297 = 641 \times 6700417$ is composite. Only five Fermat primes are known, so only these five regular polygons with odd number of sides are known to be constructible.

\begin{theorem}[Constructibility Criterion]\index{constructibility criterion}
A regular $n$-gon is constructible if and only if the Galois group $\operatorname{Gal}(\Q(\zeta_n)/\Q) \cong (\Z/n\Z)^*$ is a 2-group, i.e., has order $2^m$ for some $m$.
\end{theorem}

\section{The 17-gon Construction}

Gauss used Gaussian periods to express $\cos(2\pi/17)$ in nested square roots \index{nested radical}. Let $\zeta = e^{2\pi i/17}$. The periods of length 2 are:
\begin{align*}
\eta_1 &= \zeta + \zeta^9 + \zeta^{13} + \zeta^{15} + \zeta^{16} + \zeta^8 + \zeta^4 + \zeta^2 \\
\eta_2 &= \zeta^3 + \zeta^{10} + \zeta^5 + \zeta^{11} + \zeta^{14} + \zeta^7 + \zeta^{12} + \zeta^6
\end{align*}

These satisfy $\eta_1 + \eta_2 = -1$, $\eta_1 \eta_2 = -4$. Then $\eta_1 = \frac{-1 + \sqrt{17}}{2}$, $\eta_2 = \frac{-1 - \sqrt{17}}{2}$.

\begin{theorem}[Gaussian Periods for 17-gon]
The Gaussian periods of length 4 are:
\begin{align*}
\eta_1^{(1)} &= \zeta + \zeta^4 + \zeta^{13} + \zeta^{16} = \frac{-1 + \sqrt{17}}{4} + \frac{\sqrt{34-2\sqrt{17}}}{4} \\
\eta_1^{(2)} &= \zeta^2 + \zeta^8 + \zeta^9 + \zeta^{15} = \frac{-1 + \sqrt{17}}{4} - \frac{\sqrt{34-2\sqrt{17}}}{4} \\
\eta_2^{(1)} &= \zeta^3 + \zeta^5 + \zeta^{12} + \zeta^{14} = \frac{-1 - \sqrt{17}}{4} + \frac{\sqrt{34+2\sqrt{17}}}{4} \\
\eta_2^{(2)} &= \zeta^6 + \zeta^7 + \zeta^{10} + \zeta^{11} = \frac{-1 - \sqrt{17}}{4} - \frac{\sqrt{34+2\sqrt{17}}}{4}
\end{align*}
\end{theorem}

Continuing this process yields the nested square root expression for $\cos(2\pi/17)$:
\[
\cos\frac{2\pi}{17} = \frac{-1 + \sqrt{17} + \sqrt{34-2\sqrt{17}} + 2\sqrt{17+3\sqrt{17}-\sqrt{34-2\sqrt{17}}-2\sqrt{34+2\sqrt{17}}}}{16}.
\]

The construction requires four nested square roots, corresponding to the four bisections of the period equations. This is consistent with the fact that $17 = 2^{2^2} + 1$ is a Fermat prime.

\section{Python Implementation}

In \texttt{python/heptadecagon.py}:

\begin{lstlisting}[style=python, caption={Key functions in python/heptadecagon.tex}]
cyclotomic_poly(n)              # Phi_n(x) coefficients
constructible_ngons(limit)      # List constructible n
heptadecagon_vertices()         # Coordinates of 17-gon
gaussian_periods(n, length)     # Gaussian periods
cos_2pi_over_17()              # Exact expression
construct_polygon(n)            # General construction
\end{lstlisting}

\section{Numerical Examples}

\begin{example}[17-gon Vertices]
\begin{lstlisting}[style=python]
>>> from heptadecagon import heptadecagon_vertices
>>> verts = heptadecagon_vertices()
>>> for i, (x, y) in enumerate(verts[:4]):
...     print(f"V{i}: ({x:.6f}, {y:.6f})")
V0: (1.000000, 0.000000)
V1: (0.932472, 0.361242)
V2: (0.739009, 0.673696)
V3: (0.445738, 0.895163)
\end{lstlisting}
\end{example}

\begin{example}[Constructible $n$-gons]
\begin{lstlisting}[style=python]
>>> from heptadecagon import constructible_ngons
>>> constructible_ngons(50)
[3, 4, 5, 6, 8, 10, 12, 15, 16, 17, 20, 24, 30, 32, 34, 40, 48]
\end{lstlisting}
\end{example}

\begin{example}[Exact $\cos(2\pi/17)$]
The exact value from nested square roots:
\begin{lstlisting}[style=python]
>>> from heptadecagon import cos_2pi_over_17
>>> import math
>>> cos_val = cos_2pi_over_17()
>>> print(f"cos(2π/17) = {cos_val:.15f}")
>>> print(f"Numerical:   {math.cos(2*math.pi/17):.15f}")
cos(2π/17) = 0.932472229404356
Numerical:   0.932472229404356
\end{lstlisting}
The nested radical expression is exact to machine precision.
\end{example}

\section{Exercises}

\begin{exercise}
Derive the nested square root expression for $\cos(2\pi/17)$ step by step using Gaussian periods.
\end{exercise}

\begin{exercise}
Implement the construction of the 257-gon using Gaussian periods. How many square root nestings are needed?
\end{exercise}

\begin{exercise}
Prove that $\Phi_n(x)$ is irreducible over $\Q$ using Eisenstein's criterion \index{Eisenstein's criterion} on $\Phi_p(x+1)$ for prime $p$.
\end{exercise}

\begin{exercise}
Show that the coordinates of the 17-gon vertices lie in the field $\Q(\sqrt{17}, \sqrt{34-2\sqrt{17}}, \ldots)$.
\end{exercise}

\begin{exercise}
Compute the minimal polynomial of $\cos(2\pi/17)$ over $\Q$ and verify it has degree 8.
\end{exercise}

\begin{exercise}
Generalize Gauss's construction to prove that a regular polygon with $n$ sides is constructible iff $n$ is a product of a power of 2 and distinct Fermat primes.
\end{exercise}



\section{Exercise Solutions}

\begin{solution}
The regular heptadecagon (17-gon) has interior angle $\frac{(17-2) \cdot 180}{17} = \frac{2700}{17} \approx 158.82$. The central angle is $360/17 \approx 21.18$. Using Gauss's construction, the vertices are at angles $2\pi k/17$ for $k = 0, \ldots, 16$, with coordinates given by nested square roots.
\end{solution}

\begin{solution}
The nested square root expression for $\cos(2\pi/17)$ is derived by successive bisection of periods. Starting from $\eta_1 = \frac{-1+\sqrt{17}}{2}$, we solve for sub-periods of length 4, then 2, then 1. Each bisection introduces one additional square root. The final expression involves four nested square roots, consistent with $\varphi(17)/2 = 8$ and $8 = 2^3$.
\end{solution}

\begin{solution}
The minimal polynomial of $\cos(2\pi/17)$ is found by expressing $2\cos(2\pi/17) = \zeta + \zeta^{-1}$ and using the cyclotomic polynomial $\Phi_{17}(x) = x^{16} + x^{15} + \cdots + x + 1$. Substituting $x = e^{i\theta}$ and using $\cos\theta = \frac{e^{i\theta}+e^{-i\theta}}{2}$, we derive an 8th-degree polynomial with integer coefficients.
\end{solution}

\begin{solution}
Gauss's criterion for constructibility extends to all $n$: a regular $n$-gon is constructible iff $\varphi(n)$ is a power of 2. Since $\varphi(n) = \prod \varphi(p_i^{e_i})$, we need $\varphi(p^e) = p^{e-1}(p-1)$ to be a power of 2. This requires $e = 1$ (so $p^{e-1}=1$) and $p-1 = 2^k$ (so $p = 2^k+1$ is a Fermat prime). Thus $n = 2^m \cdot F_{i_1} \cdots F_{i_r}$.
\end{solution}


\section{References}

\begin{itemize}
    \item \cite{gauss-disq} -- \textit{Disquisitiones Arithmeticae}, Articles 355--366
    \item \cite{cox-galois} -- Cox, \textit{Galois Theory}, Ch. 11
    \item \cite{richmond-construction} -- Richmond, "A Construction for a Regular Polygon of Seventeen Sides" (1893)
    \item \cite{de-temple} -- DeTemple, "Carlyle Circles and the Lemoine Simplicity of Polygon Constructions" (1991)
    \item \cite{wantzel} -- Wantzel (1837), proof of necessity
\end{itemize}